\section{Phát biểu bài toán}
Trong thế giới công nghệ không ngừng phát triển, khái niệm về một trang web thương mại điện tử (e-commerce) sử dụng hệ thống gợi ý (recommender system) để cá nhân hóa trải nghiệm người dùng không còn xa lạ, đặc biệt là trong lĩnh vực thiết bị điện tử và máy tính. Mặc dù các nền tảng thương mại điện tử đã trở thành kênh tiếp thị và bán hàng tiêu chuẩn cho các doanh nghiệp này, vẫn tồn tại nhu cầu cải thiện các phương thức tương tác người dùng hiện tại, cụ thể là hỗ trợ người dùng truy xuất hiệu quả các sản phẩm mong muốn. Trong khi các trang web này cung cấp đa dạng sản phẩm, giao diện người dùng và hệ thống điều hướng hiện tại lại đặt ra những thách thức cho khách hàng trong việc tìm kiếm nhanh chóng và hiệu quả các máy tính và thiết bị cụ thể, dẫn đến trải nghiệm người dùng chưa tối ưu.
\newline
\newline
Thiết kế tương tác người dùng hiện tại của các trang web thương mại điện tử trong lĩnh vực máy tính và thiết bị điện tử còn thiếu các tính năng và chức năng cần thiết để nâng cao trải nghiệm duyệt và tìm kiếm. Người dùng thường gặp khó khăn trong việc định vị sản phẩm mong muốn do bộ lọc tìm kiếm hạn chế, phân loại không đầy đủ và thiếu các tùy chọn điều hướng trực quan. Điều này dẫn đến sự thất vọng, gia tăng thời gian duyệt web và nguy cơ mất hứng thú trong việc khám phá toàn bộ các sản phẩm được cung cấp.
\newline
\newline
Một trong những thách thức cụ thể nằm ở việc sử dụng các bộ lọc để thu hẹp kết quả tìm kiếm, điều này thường đòi hỏi người dùng phải có một mức độ kiến thức miền (domain knowledge) nhất định về thiết bị, điện tử hoặc máy tính. Điều này có thể gây nhầm lẫn cho những người dùng mới tiếp cận lĩnh vực này và có thể không quen thuộc với các thuật ngữ cụ thể hoặc thông số kỹ thuật. Tương tự, việc tìm kiếm sử dụng tìm kiếm toàn văn (full-text search) cũng gây khó khăn, vì người dùng cần phải hiểu rõ cách xây dựng các truy vấn tìm kiếm hiệu quả để truy xuất các sản phẩm đáp ứng nhu cầu của họ.
\newline
\newline
Tóm lại, bằng cách trao quyền cho người dùng tìm kiếm các sản phẩm đáp ứng nhu cầu mà không yêu cầu chuyên môn kỹ thuật sâu rộng hoặc phải chật vật diễn đạt nhu cầu bằng từ ngữ chính xác, các trang web có thể thu hút cơ sở khách hàng rộng lớn hơn, thúc đẩy tỷ lệ chuyển đổi và thiết lập lợi thế cạnh tranh trên thị trường.